\documentclass[12pt, a4paper]{article}

% ============================================================
%  PACOTES ESSENCIAIS
% ============================================================
\usepackage[utf8]{inputenc}
\usepackage[T1]{fontenc}
\usepackage[brazil]{babel}
\usepackage[top=3cm, bottom=2cm, left=3cm, right=2cm]{geometry}
\usepackage{graphicx}
\usepackage{booktabs}
\usepackage[hidelinks]{hyperref}
\usepackage{amsmath}
\usepackage{setspace}
\usepackage{indentfirst}
\usepackage{cite}
\usepackage{array}
\usepackage{multirow}
\usepackage{float}
\usepackage{url}
\usepackage{microtype}
\usepackage{tikz}
\usetikzlibrary{shapes.geometric, arrows.meta, positioning}

\onehalfspacing
\setlength{\parindent}{1.25cm}

% ============================================================
%  INÍCIO DO DOCUMENTO
% ============================================================
\begin{document}

% ============================================================
%  CABEÇALHO: TÍTULO, AUTORIA E FILIAÇÃO
% ============================================================
\begin{center}
{\large\bfseries
Desenvolvimento e Validação de um Sistema Automatizado para Detecção de Estratégias \textit{Call2Go} em Vídeos do YouTube: Uma Análise de Confiabilidade Metodológica no Mercado Musical Digital Brasileiro}

\vspace{0.9cm}

{\normalsize
% INSTRUÇÃO: substitua os campos entre colchetes pelos dados reais
\textbf{[Nome Completo do Autor]}\\
Graduando em [Curso de Graduação]\\
[Nome da Instituição de Ensino Superior]\\
\texttt{[email.autor@instituicao.br]}}

\vspace{0.5cm}

{\normalsize
\textbf{Orientador: [Nome Completo do Orientador]}\\
[Departamento / Curso], [Nome da Instituição]\\
\texttt{[email.orientador@instituicao.br]}}

\vspace{0.4cm}

{\small [Rua/Av. Nome, Número, Cidade --- Estado, CEP XXXXX-XXX, Brasil]}

\vspace{0.2cm}
{\small \today}
\end{center}

\vspace{0.4cm}
\hrule
\vspace{0.5cm}

% ============================================================
%  RESUMO (PORTUGUÊS)
% ============================================================
\renewcommand{\abstractname}{Resumo}
\begin{abstract}
O crescimento acelerado do mercado musical digital brasileiro elevou plataformas como YouTube e Spotify à condição de infraestrutura crítica para a circulação, descoberta e monetização de conteúdo musical. Nesse contexto, emergiu a prática denominada \textit{Call2Go} --- a inserção de chamadas de ação nas descrições de vídeos do YouTube com direcionamento explícito ao Spotify ---, adotada por artistas e \textit{labels} como estratégia de conversão de audiência entre plataformas, ainda que sem respaldo empírico consolidado sobre sua eficácia ou sobre a confiabilidade de ferramentas automatizadas aptas a detectá-la em larga escala. Este trabalho apresenta o desenvolvimento e a validação de um sistema automatizado baseado em expressões regulares (\textit{regex}) multicamada para detecção de estratégias Call2Go, aplicado a um corpus de 980 vídeos dos 49 artistas brasileiros mais visualizados, coletados via YouTube Data API v3 e Spotify Web API. O protocolo metodológico adota um ciclo de validação cruzada que contrasta, em três níveis analíticos (descrição do vídeo, perfil do canal e abordagem combinada), a classificação automatizada com anotação humana de referência (\textit{ground truth}) construída sobre uma amostra aleatória estratificada de 50 vídeos. Os resultados revelam que, no nível combinado, o sistema alcança acurácia, precisão e \textit{recall} de 100\%, validando a abordagem como instrumento metodológico confiável para análise de conteúdo em escala. Adicionalmente, o teste Mann-Whitney U não encontrou evidências estatisticamente significativas de que a adoção de Call2Go impacta o volume de visualizações no YouTube ($U = 118.004$; $p = 0{,}360$) nem a popularidade no Spotify ($U = 169.415{,}5$; $p = 0{,}998$). A análise bidirecional de correlação de Spearman identificou uma relação unidirecional no sentido Spotify $\to$ YouTube ($\rho = 0{,}392$; $p = 0{,}005$), sugerindo que a proeminência em plataformas de streaming correlaciona-se positivamente com o engajamento no YouTube, mas não o contrário. O estudo contribui com uma metodologia reprodutível, documentada e validada para a análise automatizada de práticas de \textit{cross-platform marketing} musical.
\end{abstract}

\noindent\textbf{Palavras-chave:} Call2Go; detecção automatizada de conteúdo; expressões regulares; validação cruzada; \textit{cross-platform marketing} musical.

\vspace{0.6cm}
\hrule
\vspace{0.5cm}

% ============================================================
%  ABSTRACT (INGLÊS)
% ============================================================
\section*{Abstract}

The rapid expansion of the Brazilian digital music market has elevated platforms such as YouTube and Spotify to the role of critical infrastructure for the circulation, discovery, and monetization of musical content. In this setting, the practice known as \textit{Call2Go} --- the insertion of call-to-action elements in YouTube video descriptions directing viewers toward Spotify --- has been adopted by artists and labels as an audience conversion strategy, yet without consolidated empirical support regarding its effectiveness or the reliability of automated tools designed to detect it at scale. This paper presents the development and validation of a multilayer regex-based automated system for Call2Go detection, applied to a corpus of 980 videos from the 49 most-viewed Brazilian artists, collected via the YouTube Data API v3 and the Spotify Web API. The methodological protocol adopts a cross-validation cycle contrasting, at three analytical levels (video description, channel profile, and combined approach), automated classification against a human reference annotation (\textit{ground truth}) built on a stratified random sample of 50 videos. Results show that at the combined level the system achieves 100\% accuracy, precision, and recall, validating the approach as a reliable methodological instrument for content analysis at scale. Furthermore, the Mann-Whitney U test found no statistically significant evidence that Call2Go adoption impacts YouTube view volume ($U = 118{,}004$; $p = 0.360$) or Spotify popularity ($U = 169{,}415.5$; $p = 0.998$). Bidirectional Spearman correlation analysis identified a unidirectional relationship in the Spotify $\to$ YouTube direction ($\rho = 0.392$; $p = 0.005$), suggesting that streaming platform prominence positively correlates with YouTube engagement, but not the reverse. The study contributes a reproducible, documented, and validated methodology for the automated analysis of cross-platform marketing practices in the music industry.

\noindent\textbf{Keywords:} Call2Go; automated content detection; regular expressions; cross-validation; digital music cross-platform marketing.

\vspace{0.4cm}
\hrule
\vspace{0.4cm}

% ============================================================
%  1. INTRODUÇÃO
% ============================================================
\section{Introdução}

A transformação digital da indústria fonográfica nas últimas duas décadas reconfigurou profundamente as relações entre produção, distribuição e consumo de música. O relatório anual da Federação Internacional da Indústria Fonográfica registrou, em 2023, que as receitas provenientes de serviços de \textit{streaming} responderam por 67\% do total global da indústria da música gravada, marcando uma inflexão histórica no modelo de negócios do setor \cite{ifpi2023}. O Brasil ocupa posição de destaque nesse cenário: é o décimo maior mercado fonográfico do planeta e exibe taxas de crescimento em \textit{streaming} que superam consistentemente a média global \cite{ifpi2023}. Dentro desse ecossistema, duas plataformas assumiram protagonismo singular: o YouTube, maior repositório de vídeos do mundo e principal canal de descoberta musical, e o Spotify, líder global em \textit{streaming} de áudio, com base de usuários ativos que ultrapassa 600 milhões. A coexistência dessas duas plataformas na trajetória de consumo do ouvinte brasileiro criou uma dinâmica de complementaridade e competição que exige das equipes de \textit{marketing} artístico estratégias específicas de gerenciamento de audiência entre plataformas.

Nesse contexto, a estratégia denominada \textit{Call2Go} emergiu como prática amplamente difundida na indústria musical doméstica. O termo descreve qualquer elemento inserido na descrição de um vídeo do YouTube --- seja um hiperlink direto para o Spotify, seja uma chamada textual de ação (\textit{call-to-action}) --- com o propósito explícito de redirecionar o espectador para a plataforma de \textit{streaming} de áudio. A lógica subjacente à estratégia fundamenta-se na hipótese de que a audiência conquistada no YouTube, plataforma de acesso gratuito e onipresente, pode ser convertida em assinantes ou ouvintes recorrentes no Spotify, cuja monetização por \textit{stream} é diretamente mais vantajosa para artistas e distribuidoras. A viabilidade técnica da estratégia é imediata: o YouTube permite o preenchimento irrestrito das descrições de vídeo, tornando trivial a inserção de links e textos de redirecionamento. Seu impacto efetivo, contudo, permanece empiricamente indeterminado \cite{wlomert2016}.

A investigação científica sobre as dinâmicas de \textit{cross-platform} na música digital ainda se encontra em estágio formativo \cite{aguiar2018}. Os estudos existentes concentram-se predominantemente em análises agregadas de mercado ou em fenômenos observados em plataformas isoladas, sem abordar de forma sistemática os mecanismos de transferência de audiência entre YouTube e Spotify \cite{napoli2012}. Ademais, a ausência de ferramentas automatizadas validadas para a detecção de Call2Go em larga escala representa um obstáculo metodológico significativo: a anotação manual de grandes corpora é ao mesmo tempo impraticável operacionalmente e sujeita a vieses de subjetividade entre anotadores. A lacuna identificada é, portanto, dupla: empírica, pela ausência de evidências sobre a eficácia da estratégia, e metodológica, pela inexistência de instrumentos confiáveis para sua detecção automatizada em escala.

Este trabalho endereça diretamente essa lacuna ao propor, implementar e validar um sistema automatizado de detecção de Call2Go baseado em expressões regulares multicamada. O foco central da investigação recai sobre a \textbf{confiabilidade metodológica} da ferramenta desenvolvida, mensurada por meio de um protocolo de validação cruzada entre classificação automatizada e anotação humana de referência (\textit{ground truth}). Secundariamente, o trabalho emprega os dados coletados para testar hipóteses sobre o impacto estatístico do Call2Go no engajamento em ambas as plataformas e para investigar a direcionalidade da correlação entre métricas do YouTube e do Spotify.

O artigo está organizado da seguinte forma: a Seção~\ref{sec:problema} delimita o problema de pesquisa e apresenta as hipóteses formalizadas; a Seção~\ref{sec:literatura} revisa a literatura relevante nas três dimensões temáticas do trabalho; a Seção~\ref{sec:objetivos} expõe os objetivos geral e específicos; a Seção~\ref{sec:metodologia} detalha o pipeline metodológico, incluindo coleta de dados, sistema de detecção e protocolo de validação; a Seção~\ref{sec:resultados} apresenta e discute os resultados obtidos; e a Seção~\ref{sec:conclusao} encerra com as conclusões, limitações e direcionamentos futuros.

% ============================================================
%  2. PROBLEMA DE PESQUISA
% ============================================================
\section{Problema de Pesquisa}
\label{sec:problema}

O problema central deste trabalho pode ser formulado em dois planos interdependentes. No plano instrumental, questiona-se: \textit{é possível construir um sistema automatizado baseado em expressões regulares capaz de detectar estratégias Call2Go em descrições de vídeos do YouTube com confiabilidade estatisticamente equivalente à anotação humana?} No plano substantivo, questiona-se: \textit{o uso da estratégia Call2Go por artistas brasileiros no YouTube está associado a diferenças mensuráveis no volume de visualizações na mesma plataforma ou na popularidade no Spotify?}

A ambiguidade característica dos textos de descrição de vídeo no YouTube --- repletos de abreviações, emojis, URLs encurtadas, redirecionadores de terceiros e menções narrativas não-intencionais a plataformas --- torna a detecção automatizada um problema não-trivial. Um sistema baseado em regras simples incorre em alta taxa de falsos positivos ao tratar qualquer menção textual à palavra ``Spotify'' como indicativo de Call2Go, sem distinguir entre uma chamada de ação intencional e uma referência contextual (por exemplo, ``música do artista X no topo dos charts do Spotify''). Ao mesmo tempo, a dependência exclusiva de links diretos subestima a detecção, pois parte dos links de Call2Go é veiculada por meio de serviços de redirecionamento (\textit{smart links}) como \texttt{lnk.to}, \texttt{bit.ly} ou \texttt{smarturl.it}, cujo domínio de destino final (Spotify) não é imediatamente visível no texto da descrição.

A delimitação do problema exige, portanto, a operacionalização formal de três dimensões analíticas: (i) a \textbf{taxonomia de detecção}, definindo as categorias mutuamente exclusivas e coletivamente exaustivas em que se enquadra cada vídeo (\texttt{link\_direto}, \texttt{texto\_implicito} ou \texttt{nenhum}); (ii) a \textbf{hierarquia de fontes}, reconhecendo que a evidência de Call2Go pode estar no vídeo individual, no perfil do canal ou em ambos, com implicações distintas de confiabilidade; e (iii) o \textbf{escopo da validação}, determinando o tamanho e a estratégia de amostragem para a construção do \textit{ground truth} humano. A ausência de qualquer uma dessas dimensões na literatura prévia específica ao mercado musical brasileiro justifica a originalidade metodológica da proposta.

As hipóteses de pesquisa formalizadas são:

\begin{enumerate}
    \item[$H_0$] (Hipótese nula da validação): O detector automatizado baseado em \textit{regex} não apresenta concordância estatisticamente significativa com a classificação humana, i.e., a acurácia do sistema é inferior a 85\%.
    \item[$H_1$] (Hipótese alternativa da validação): O detector automatizado atinge acurácia $\geq 85\%$ em relação ao \textit{ground truth} humano, validando a abordagem como metodologicamente confiável.
    \item[$H_2$] (Hipótese secundária --- YouTube): Vídeos com Call2Go geram, em média, volume significativamente maior de visualizações no YouTube em comparação com vídeos sem a estratégia.
    \item[$H_3$] (Hipótese secundária --- cross-platform): A presença de Call2Go nos vídeos do YouTube de um artista está positivamente associada a maior pontuação de popularidade do artista no Spotify.
\end{enumerate}

% ============================================================
%  3. REVISÃO DA LITERATURA
% ============================================================
\section{Revisão da Literatura}
\label{sec:literatura}

\subsection{Marketing Digital na Indústria Musical Contemporânea}

A digitalização da cadeia produtiva musical transformou não apenas os canais de distribuição, mas os próprios modelos de relacionamento entre artistas, intermediários e audiências. Aguiar e Waldfogel \cite{aguiar2018} documentaram empiricamente que a entrada de plataformas de distribuição digital alterou as dinâmicas de promoção da música, deslocando poder das grandes gravadoras para distribuidoras independentes e, em determinados contextos, para os próprios artistas. Nesse novo ecossistema, a capacidade de gestão estratégica da presença multiplataforma tornou-se competência central para artistas que buscam maximizar alcance e monetização. O YouTube e o Spotify, embora frequentemente tratados como concorrentes na disputa pela atenção do ouvinte, exercem funções complementares na jornada de consumo musical: o YouTube opera primariamente como plataforma de descoberta e videoconsumo, ao passo que o Spotify consolida o consumo recorrente em forma de \textit{stream} contabilizável para fins de remuneração \cite{wlomert2016}.

As estratégias de \textit{cross-platform marketing} na música digital seguem uma lógica de funil: o conteúdo audiovisual no YouTube atua como ponto de entrada de alta viralidade e baixo custo de aquisição, enquanto as plataformas de \textit{streaming} de áudio representam o destino de conversão de valor. Nesse modelo, o Call2Go funciona como o mecanismo de transição --- o elo operacional que transforma um espectador passivo em um ouvinte ativo e mensurável. A eficácia empírica dessa conversão, contudo, permanece pouco estudada. Wlömert e Papies \cite{wlomert2016} demonstraram que a entrada do Spotify no mercado gerou efeitos assimétricos sobre as receitas da indústria, com impactos distintos sobre as vendas digitais e sobre o acesso gratuito a conteúdo, mas sem abordar especificamente os mecanismos de direcionamento interplataforma originados no YouTube.

O mercado brasileiro apresenta características estruturais que potencializam a relevância da questão. A penetração do YouTube no Brasil é das mais altas do mundo em termos relativos, com o país figurando consistentemente entre os maiores mercados da plataforma por volume de horas assistidas \cite{ifpi2023}. Simultaneamente, o Brasil experimenta uma das transições mais rápidas do mundo de consumo musical gratuito para streaming pago, com artistas do segmento sertanejo, pagode e funk acumulando bilhões de visualizações no YouTube enquanto buscam migrar sua base para plataformas remuneradas. Essa configuração torna o mercado musical brasileiro um laboratório privilegiado para a investigação das dinâmicas de \textit{cross-platform} na música popular de massa.

A ausência de literatura científica brasileira específica sobre Call2Go como estratégia de \textit{marketing} artístico representa tanto a principal lacuna que este trabalho busca endereçar quanto uma oportunidade de contribuição original para a área. Os trabalhos internacionais mais próximos focam na análise de estratégias de CTA (\textit{call-to-action}) em contextos de \textit{e-commerce} e redes sociais genéricas \cite{napoli2012}, sem contemplar as especificidades do conteúdo musical e das audiências de entretenimento.

\subsection{Análise Automatizada de Conteúdo Textual em Plataformas Digitais}

A análise de conteúdo como método científico possui raízes que remontam ao século XX e foi sistematizada de forma seminal por Krippendorff \cite{krippendorff2004}, que a define como ``uma técnica de pesquisa para fazer inferências replicáveis e válidas a partir de textos (ou outros materiais significativos) para os contextos de seu uso''. A transição do método para ambientes digitais ampliou dramaticamente a escala dos corpora analisáveis, tornando a automação não apenas desejável, mas necessária. Plataformas como YouTube e Twitter geram diariamente volumes de texto que tornam a anotação manual integralmente impraticável para fins de pesquisa quantitativa em larga escala.

No campo específico da recuperação de informação musical (\textit{Music Information Retrieval} --- MIR), Schedl, Gómez e Urbano \cite{schedl2014} catalogaram as principais abordagens de análise automatizada de metadados musicais, incluindo o processamento de texto associado a conteúdo audiovisual. As abordagens de detecção de padrões textuais em descrições de vídeo variam desde métodos baseados em regras simples (\textit{keyword matching}) até modelos de aprendizado de máquina supervisionados, passando por sistemas híbridos. Os métodos baseados em regras, como o sistema \textit{regex} proposto neste trabalho, caracterizam-se por alta interpretabilidade, custo computacional negligenciável e comportamento determinístico --- propriedades especialmente valiosas em contexto acadêmico, onde a reprodutibilidade e a auditabilidade do processo de classificação constituem requisitos metodológicos \cite{salganik2019}. A principal vulnerabilidade desses sistemas é a sensibilidade ao contexto: padrões extraídos de um corpus específico podem ter baixa generalização para domínios distintos, exigindo processo cuidadoso de construção, teste e validação do conjunto de regras.

Celma \cite{celma2010} explorou os desafios da recomendação musical automatizada, destacando que a esparsidade e o ruído inerentes aos metadados de vídeo e música dificultam a extração consistente de informações estruturadas. O fenômeno é particularmente agudo em descrições de vídeos do YouTube, onde não há padronização de formato, a presença de URLs encurtadas e redirecionadores de terceiros oculta o destino final dos links, e menções narrativas a serviços de streaming são semanticamente indistinguíveis de chamadas de ação sem análise contextual aprofundada. O sistema desenvolvido neste trabalho enfrenta diretamente esses desafios por meio de uma arquitetura de detecção em três camadas, conforme detalhado na Seção~\ref{sec:metodologia}.

A distinção entre menções narrativas e chamadas de ação constitui um dos problemas mais recorrentes em análise automatizada de descrições de vídeo. Um artista pode mencionar o Spotify em frases como ``alcançamos 10 milhões de plays no Spotify'' (menção narrativa, não-Call2Go) ou ``ouça no Spotify'' (chamada de ação, Call2Go). Sistemas baseados em \textit{keyword matching} simples tratam ambos os casos de forma idêntica, produzindo falsos positivos que comprometem a validade da análise. O desenvolvimento de padrões de exclusão específicos para menções narrativas --- como referências a rankings, charts e estatísticas --- é um dos diferenciais técnicos do sistema proposto.

\subsection{Confiabilidade em Sistemas de Classificação Automatizada}

A avaliação da confiabilidade de sistemas de classificação automatizada em relação a referências humanas constitui um problema central em pesquisa computacional de ciências sociais e humanidades digitais. O índice kappa de Cohen \cite{cohen1960} representa a medida pioneira de concordância entre dois anotadores além do nível esperado pelo acaso, tendo sido amplamente adotado como padrão metodológico em estudos de conteúdo. Fleiss \cite{fleiss1971} estendeu o índice para múltiplos anotadores, consolidando a família de métricas kappa como instrumento de referência na verificação de inter-rater reliability. Ambas as métricas, contudo, pressupõem anotações nominais categóricas e não capturam as nuances de sistemas com hierarquias de confiança variável, motivando o uso complementar de medidas como acurácia, precisão, recall e F1-score adotadas no presente trabalho.

Salganik \cite{salganik2019} argumenta que pesquisas computacionais em ciências sociais enfrentam um dilema fundamental entre escala e validade: métodos automatizados permitem analisar corpora de milhões de registros, mas sua validade depende de etapas de verificação humana que, por razões práticas, recaem sobre amostras. O design da amostra para construção do \textit{ground truth} é, portanto, uma decisão metodológica de primeira ordem. A estratégia de amostragem aleatória simples com semente fixa (seed), adotada neste trabalho, garante a reprodutibilidade da amostra e a ausência de viés de seleção, enquanto o tamanho amostral de 50 vídeos, selecionados de um universo de 980, corresponde a uma taxa de amostragem de 5,1\%, considerada adequada para estimativas de proporção com erro de 14\% ao nível de confiança de 95\%.

A literatura de aprendizado de máquina aplicado à análise de conteúdo distingue entre confiabilidade interna (consistência do sistema consigo mesmo em condições idênticas) e confiabilidade externa (concordância com referência humana). Para sistemas baseados em regras determinísticas como o proposto, a confiabilidade interna é por definição perfeita. A confiabilidade externa, medida pelo confronto com o \textit{ground truth}, é a métrica relevante para este trabalho e constitui a base para a aceitação ou rejeição de $H_1$. A Tabela~\ref{tab:literatura} apresenta uma síntese comparativa das principais abordagens relatadas na literatura para análise automatizada de conteúdo em contextos similares, evidenciando os compromissos entre escala, custo e validade em cada estratégia.

\begin{table}[H]
\centering
\caption{Comparativo de abordagens para análise automatizada de conteúdo em plataformas digitais.}
\label{tab:literatura}
\renewcommand{\arraystretch}{1.3}
\begin{tabular}{@{}lllll@{}}
\toprule
\textbf{Abordagem} & \textbf{Mecanismo} & \textbf{Escala} & \textbf{Validação} & \textbf{Limitação principal} \\
\midrule
Anotação manual & Humano & Pequena & $\checkmark$ Alta & Não escalável \\
\textit{Keyword matching} & Palavras-chave fixas & Grande & Baixa & Alta taxa de FP \\
\textit{Regex} básico & Padrões fixos & Grande & Média & Sensível ao contexto \\
ML supervisionado & Classificador treinado & Grande & Média & Necessita corpus rotulado \\
NLP profundo & \textit{Transformers} (BERT, etc.) & Grande & Alta & Alto custo computacional \\
\textbf{Este trabalho} & \textbf{\textit{Regex} multicamada} & \textbf{Grande} & \textbf{$\checkmark$ Alta (100\%)} & \textbf{Domínio-específico} \\
\bottomrule
\end{tabular}
\end{table}

% ============================================================
%  4. OBJETIVOS DA PESQUISA
% ============================================================
\section{Objetivos da Pesquisa}
\label{sec:objetivos}

\subsection{Objetivo Geral}

Desenvolver, implementar e validar um sistema automatizado baseado em expressões regulares para a detecção de estratégias Call2Go em descrições de vídeos do YouTube, aplicado ao corpus dos artistas brasileiros mais visualizados, e avaliar a confiabilidade metodológica da ferramenta por meio de um protocolo de validação cruzada entre classificação automática e anotação humana de referência.

\subsection{Objetivos Específicos}

\begin{enumerate}

    \item \textbf{Construir uma base de artistas reprodutível e verificável}, derivada de playlists oficiais do Spotify (Top 50 Brasil, Viral 50 Brasil e Top Hits Brasil), integrando métricas de popularidade e identificadores de canal do YouTube, de forma a garantir que a amostra represe os artistas com maior relevância na plataforma principal de descoberta musical no Brasil. A reprodutibilidade da fonte de dados é condição necessária para a replicabilidade da pesquisa.

    \item \textbf{Coletar e estruturar um corpus de vídeos do YouTube}, composto pelos 20 vídeos mais visualizados de cada artista, totalizando até 980 vídeos, utilizando a YouTube Data API v3 com estratégia de consulta otimizada por quotas (\texttt{playlistItems.list}), e complementado por scraping estruturado das informações de canal (aba ``Sobre''), que contêm links externos não acessíveis diretamente via API.

    \item \textbf{Desenvolver um motor de detecção Call2Go em arquitetura de três camadas analíticas} (vídeo, canal textual e canal via \textit{scraping} estruturado), capaz de distinguir entre links diretos ao Spotify, redirecionadores de terceiros com destino ao Spotify e menções textuais implícitas, excluindo automaticamente menções narrativas que não constituem chamadas de ação, por meio de padrões de exclusão formalmente definidos.

    \item \textbf{Construir um \textit{ground truth} humano} sobre amostra aleatória de 50 vídeos (seed = 42) extraída do corpus total, utilizando protocolo semi-automático de pré-anotação que maximiza a eficiência da revisão humana ao apresentar ao anotador apenas os campos de evidência relevantes, garantindo cobertura plena de confiança alta ($\geq$~98\% dos vídeos) antes da validação final.

    \item \textbf{Avaliar a confiabilidade do detector automatizado} nos três níveis analíticos (vídeo isolado, canal isolado e combinado vídeo + canal), reportando, para cada nível e cada classe, as métricas de acurácia global, precisão, recall e F1-score, além de matriz de confusão, de forma a permitir a caracterização precisa dos tipos e magnitudes dos erros de classificação.

    \item \textbf{Testar as hipóteses secundárias de impacto} ($H_2$ e $H_3$) por meio de testes estatísticos não-paramétricos (Mann-Whitney U, $\alpha = 0{,}05$), verificando se a adoção de Call2Go está associada a diferenças estatisticamente significativas no volume de visualizações no YouTube e na popularidade no Spotify, e conduzir análise bidirecional de correlação de Spearman para determinar a direcionalidade da influência cruzada entre as duas plataformas.

\end{enumerate}

% ============================================================
%  5. METODOLOGIA DE PESQUISA
% ============================================================
\section{Metodologia de Pesquisa}
\label{sec:metodologia}

A Figura~\ref{fig:pipeline} apresenta o pipeline ETL (\textit{Extract, Transform, Load}) e de validação executado integralmente neste trabalho. O pipeline é composto por 11 etapas sequenciais, das quais as quatro primeiras são responsáveis pela extração de dados brutos, a quinta pela classificação automatizada, a sexta pelo armazenamento estruturado, as etapas sete a nove pelas análises estatísticas e as etapas dez e onze pelo ciclo de validação cruzada.

\begin{figure}[H]
\centering
\begin{tikzpicture}[
    node distance=0.95cm,
    box/.style={rectangle, draw=black!70, fill=gray!10, text width=9cm,
                text centered, minimum height=0.7cm, font=\small,
                rounded corners=2pt},
    vbox/.style={rectangle, draw=blue!60, fill=blue!8, text width=9cm,
                 text centered, minimum height=0.7cm, font=\small,
                 rounded corners=2pt},
    arr/.style={-{Stealth}, thick, draw=black!60}
]
\node[box]  (s1)  {Etapa 1: Fonte Oficial de Artistas --- Playlists Spotify Top BR};
\node[box, below of=s1]  (s2)  {Etapas 2--3: Coleta via API (YouTube Data API v3 + Spotify Web API)};
\node[box, below of=s2]  (s3)  {Etapa 4: \textit{Scraping} Estruturado dos Links de Canal (Aba ``Sobre'')};
\node[box, below of=s3]  (s4)  {Etapa 5: Detecção Automatizada de Call2Go (Motor \textit{Regex} Multicamada)};
\node[box, below of=s4]  (s5)  {Etapa 6: Construção do \textit{Data Warehouse} Relacional (SQLite)};
\node[box, below of=s5]  (s6)  {Etapas 7--9: Análises Estatísticas (EDA + Mann-Whitney U + \textit{Cross-Platform})};
\node[vbox, below of=s6]  (s7)  {Etapa 10: Amostragem Aleatória para \textit{Ground Truth} (50 vídeos, seed = 42)};
\node[vbox, below of=s7]  (s8)  {Etapa 11: Validação Cruzada Humano vs. Automático (3 Níveis)};
\draw[arr] (s1) -- (s2);
\draw[arr] (s2) -- (s3);
\draw[arr] (s3) -- (s4);
\draw[arr] (s4) -- (s5);
\draw[arr] (s5) -- (s6);
\draw[arr] (s6) -- (s7);
\draw[arr] (s7) -- (s8);
\end{tikzpicture}
\caption{Pipeline ETL e de validação do sistema Call2Go Detector (etapas em cinza: processamento; etapas em azul: ciclo de validação cruzada).}
\label{fig:pipeline}
\end{figure}

\subsection{Fonte de Dados e Critérios de Seleção}

A construção reprodutível da base de artistas constitui decisão metodológica crítica. A utilização de playlists editoriais do Spotify (``Top 50 Brasil'', ``Viral 50 Brasil'' e ``Top Hits Brasil'') como fonte primária garante três propriedades essenciais: (i) autoridade curatorial, pois as playlists são mantidas pela própria plataforma; (ii) relevância cultural, pois refletem os padrões de consumo contemporâneos; e (iii) reprodutibilidade, pois os identificadores de playlist (\textit{Spotify Playlist IDs}) são permanentes e publicamente acessíveis via Spotify Web API. O processo de seleção consistiu em consultar as playlists, extrair os artistas únicos, enriquecê-los com métricas de \textit{viewCount} do YouTube para classifica-los por popularidade audiovisual, e aplicar deduplicação por \textit{artist\_id} para eliminar artistas repetidos entre playlists. Foi aplicado critério de popularidade mínima no Spotify ($\geq 60$ pontos em escala 0--100) para garantir que os artistas selecionados possuíam presença relevante nas duas plataformas simultaneamente. O resultado final foi uma base de 49 artistas brasileiros após a remoção do único artista identificado com mapeamento incorreto de canal no YouTube (Label MJ Records mapeada para o canal do artista estrangeiro Michael Jackson), totalizando 980 vídeos.

Para a coleta de dados do YouTube, foi utilizada a YouTube Data API v3 com uma estratégia de acesso otimizada para redução do consumo de cota diária. A consulta padrão por meio do endpoint \texttt{search.list} consome 100 unidades de cota por chamada, tornando a coleta de 50 artistas com 20 vídeos cada inviável dentro da cota diária de 10.000 unidades. A estratégia adotada utiliza o endpoint \texttt{playlistItems.list} (custo de 1 unidade por chamada), acessando a \textit{uploads playlist} de cada canal (identificada pelo padrão \texttt{UC} $\to$ \texttt{UU}) para recuperar até 200 vídeos, ordenando localmente por \texttt{viewCount} e selecionando os 20 mais visualizados. Essa abordagem reduziu o consumo estimado de cota em aproximadamente 89\%, de $\sim$5.150 para $\sim$550 unidades para 50 artistas.

As métricas do Spotify foram coletadas via Spotify Web API utilizando o fluxo de autenticação \textit{Client Credentials} (OAuth~2.0), sem necessidade de autorização do usuário. O \textit{snapshot} de popularidade e seguidores foi capturado em 30 de março de 2026, constituindo um ponto de medição transversal. O campo \texttt{popularity} do Spotify representa um índice algorítmico de 0 a 100, calculado internamente pela plataforma com base na recência e frequência dos \textit{streams}, não sendo transparente em relação à fórmula exata de cálculo, o que constitui uma limitação metodológica explicitada adiante.

\subsection{Scrapin Estruturado de Links de Canal}

A aba ``Sobre'' dos canais do YouTube contém links externos estruturados que \textbf{não são expostos pela YouTube Data API v3}. Essa lacuna, identificada durante a fase de análise exploratória, é metodologicamente crítica: os links do Spotify inseridos na seção pública do canal representam a principal fonte de Call2Go para 50,4\% dos vídeos do corpus. Para suprir essa lacuna, foi desenvolvido um \textit{scraper} em duas fases que acessa diretamente o HTML da página do canal: a Fase~1 extrai o objeto JSON \texttt{ytInitialData} da página principal e busca o elemento \texttt{channelExternalLinkViewModel}; a Fase~2 acessa a subpágina \texttt{/about} como \textit{fallback}. Um tratamento especial foi implementado para a decodificação de caracteres \texttt{\textbackslash u0026} (sequência de escape Unicode para ``\&'') presentes nos parâmetros de URL das páginas do YouTube, cuja ausência de decodificação resultava em URLs de redirecionamento truncadas e, consequentemente, não-classificáveis. Canais do tipo OAC (\textit{Official Artist Channel}) --- gerados automaticamente pelo YouTube por meio de fusão de canais --- foram tratados de forma especial: como não possuem links personalizados na aba ``Sobre'', o algoritmo busca automaticamente o canal oficial principal do artista via pesquisa interna no YouTube.

\subsection{Sistema de Detecção Call2Go: Motor \textit{Regex} Multicamada}

O núcleo técnico do trabalho é o motor de detecção Call2Go, implementado como pipeline de classificação baseado em expressões regulares em três camadas successivas de análise:

\textbf{Camada 1 --- Links diretos ao Spotify:} Detecta a presença de hiperlinks cujo domínio pertence ao Spotify ou a serviços de redirecionamento conhecidos por operar em nome do Spotify. Os padrões implementados cobrem três subtipos:
\begin{itemize}
    \item Domínios proprietários do Spotify: \texttt{open.spotify.com}, \texttt{spoti.fi} e \texttt{sptfy.com};
    \item Redirecionadores de terceiros com rótulo explícito (\textit{labeled redirect}): expressões regulares que detectam a coexistência de um domínio de redirecionamento (e.g., \texttt{lnk.to}, \texttt{smarturl.it}, \texttt{bit.ly}) com o rótulo textual ``spotify'' a uma distância máxima de 50 caracteres na sequência do link, operacionalizando o conceito de ``smart link rotulado'';
    \item Parâmetro de caminho URL com \textit{spotify}/\textit{sptfy} embutido (e.g., \texttt{/spotify} como componente de path).
\end{itemize}

\textbf{Camada 2 --- Texto implícito (chamada verbal ao Spotify):} Detecta padrões textuais de CTA sem link, cobrindo quatro construções em português:
\begin{itemize}
    \item \texttt{ou[çc]a} $+$ \texttt{.{0,50}} $+$ \texttt{spotify} (e variações com palavras intermediárias);
    \item \texttt{dispon[ií]vel no spotify};
    \item \texttt{stream.{0,50}spotify} (alcance de 50 caracteres, em substituição ao operador guloso \texttt{.*} original, que gerava falsos positivos em bios extensas);
    \item \texttt{ouvir.{0,50}spotify}.
\end{itemize}

\textbf{Filtragem de menções narrativas:} Antes da aplicação das camadas de detecção, o texto é submetido a um pré-filtro que invalida segmentos correspondentes a menções narrativas formais, definidas pelo conjunto de padrões:
\begin{itemize}
    \item \texttt{chart\textbackslash w*\textbackslash s+(?:do|no)\textbackslash s+spotify} (``charts do Spotify'');
    \item \texttt{ranking\textbackslash s+(?:do|no)\textbackslash s+spotify};
    \item \texttt{top\textbackslash s+\textbackslash d+\textbackslash s+(?:do|no)\textbackslash s+spotify};
    \item \texttt{milh[ãõ]\textbackslash w*\textbackslash s+(?:de\textbackslash s+)?(?:plays?|streams?)\textbackslash s+(?:no|do)\textbackslash s+spotify}.
\end{itemize}

A hierarquia de fontes reconhece que a evidência de Call2Go pode ser encontrada no texto da descrição do vídeo individual (\textsl{fonte:~vídeo}), no perfil textual do canal (\textsl{fonte:~canal}), nos links estruturados scrapeados da aba ``Sobre'' (\textsl{fonte:~canal/scraped}) ou em combinação de fontes (\textsl{fonte:~ambos}). A classificação final do vídeo agrega as evidências de todas as fontes disponíveis.

\subsection{Protocolo de Validação Cruzada}

O protocolo de validação cruzada foi projetado para medir a concordância entre o sistema automatizado e a classificação humana em três níveis analíticos:

\begin{itemize}
    \item \textbf{Nível 1 --- Só vídeo:} Confronta a classificação humana com a detecção automática restrita à descrição do vídeo. Mede a confiabilidade do sistema em cenário de ausência de informações de canal.
    \item \textbf{Nível 2 --- Combinado (vídeo + canal):} Confronta com a detecção automática completa, que incorpora todas as fontes disponíveis. Representa o modo de operação pleno do sistema.
    \item \textbf{Nível 3 --- Canal isolado:} Confronta a classificação humana com a detecção baseada exclusivamente no perfil do canal. Isola a contribuição do \textit{scraping} de canal para a acurácia total.
\end{itemize}

O \textit{ground truth} foi construído da seguinte forma: (i) seleção aleatória de 50 vídeos do corpus total de 980, com seed = 42, garantindo reprodutibilidade; (ii) pré-preenchimento automatizado por meio de \texttt{ground\_truth\_helper.py}, que cruza os 50 vídeos com os dados brutos completos, aplicando o detector e extraindo evidências textuais relevantes para cada caso, classificando cada vídeo com nível de confiança (ALTA / MÉDIA / BAIXA) com base na disponibilidade e consistência das evidências; (iii) revisão humana individual de cada caso, com especial atenção aos casos de confiança BAIXA e MÉDIA; (iv) resolução manual de ambiguidades (e.g., menções narrativas vs. CTAs, mismatches de channel\_id). O resultado final foi uma amostra com 49/50 vídeos de confiança ALTA e 1/50 de confiança MÉDIA (revisado e confirmado), com distribuição de 29 vídeos \texttt{link\_direto} e 21 vídeos \texttt{nenhum}.

Para os testes de hipótese $H_2$ e $H_3$, o teste de Mann-Whitney U foi selecionado por ser não-paramétrico (sem pressuposição de normalidade da distribuição de \textit{view\_count}), adequado para amostras de tamanhos distintos e robusto a outliers \cite{mann1947}. O nível de significância adotado foi $\alpha = 0{,}05$.

Para a análise bidirecional, a correlação de Spearman foi adotada por sua capacidade de detectar relações monotônicas não-lineares entre variáveis ordinais ou contínuas com distribuição não-normal, o que é especialmente relevante para métricas de plataformas digitais, notoriamente assimétricas. A análise foi conduzida em dois níveis: (i) por artista (N = 49), correlacionando métricas agregadas do YouTube (\textit{avg\_views}, \textit{avg\_engagement}, \textit{call2go\_rate}) com métricas do Spotify (\textit{popularity}, \textit{followers}); e (ii) por vídeo (N = 980), correlacionando a popularidade do artista no Spotify com o \textit{view\_count} de cada vídeo individual. O limiar de significância para a análise bidirecional foi ajustado para $\alpha = 0{,}10$, reconhecendo o tamanho amostral limitado no nível por-artista.

% ============================================================
%  6. IMPLEMENTAÇÕES, TESTES E RESULTADOS PRELIMINARES
% ============================================================
\section{Implementações, Testes e Resultados Preliminares}
\label{sec:resultados}

\subsection{Distribuição das Detecções de Call2Go}

A execução do motor de detecção sobre o corpus completo de 980 vídeos produziu a distribuição apresentada na Tabela~\ref{tab:deteccao}. O resultado central é a \textbf{ausência completa da categoria \texttt{texto\_implicito}} no corpus, que foi eliminada durante o processo de refinamento do sistema (Fase~6) após a identificação de falsos positivos sistemáticos causados pelo operador guloso \texttt{.*} nos padrões de texto implícito. Na versão corrigida com o operador de alcance limitado \texttt{.{0,50}}, todos os casos anteriormente classificados como \texttt{texto\_implicito} foram reclassificados como \texttt{nenhum}, indicando que eram menções narrativas capturadas indevidamente pelo padrão guloso.

\begin{table}[H]
\centering
\caption{Distribuição das detecções de Call2Go no corpus de 980 vídeos dos 49 artistas brasileiros.}
\label{tab:deteccao}
\renewcommand{\arraystretch}{1.3}
\begin{tabular}{@{}llrrc@{}}
\toprule
\textbf{Classificação} & \textbf{Fonte principal} & \textbf{Vídeos} & \textbf{Percentual} & \textbf{Artistas afetados} \\
\midrule
\texttt{link\_direto} & Canal (\textit{scraped}) & 494 & 50,4\% & -- \\
\texttt{link\_direto} & Ambos (vídeo + canal) & 66 & 6,7\% & -- \\
\texttt{link\_direto} & Vídeo (descrição) & 15 & 1,5\% & -- \\
\textbf{Subtotal \texttt{link\_direto}} & & \textbf{575} & \textbf{58,7\%} & \textbf{29 de 49} \\
\midrule
\texttt{texto\_implicito} & --- & 0 & 0,0\% & -- \\
\midrule
\textbf{\texttt{nenhum}} & --- & \textbf{405} & \textbf{41,3\%} & \textbf{20 de 49} \\
\midrule
\textbf{Total} & & \textbf{980} & \textbf{100,0\%} & \textbf{49} \\
\bottomrule
\end{tabular}
\end{table}

A decomposição por fonte revela que 50,4\% dos vídeos do corpus têm o Call2Go detectado exclusivamente por meio das informações de canal (aba ``Sobre'' scrapeada), o que demonstra a imprescindibilidade do \textit{scraping} estruturado que supre a lacuna da API. Apenas 15 vídeos (1,5\%) têm o Call2Go detectado exclusivamente na descrição do vídeo individual, confirmando que a prática dominante é inserir o link no perfil do canal, não em cada vídeo separadamente. Os 66 vídeos classificados como \texttt{ambos} correspondem a artistas que adotam a prática nos dois níveis simultaneamente.

\subsection{Análise Descritiva (EDA)}

A análise exploratória confrontou as distribuições de \textit{view\_count} para os grupos \texttt{link\_direto} ($n = 717$ registros no cruzamento com Spotify, após \textit{inner join}) e \texttt{nenhum} ($n = 523$ registros). Os vídeos com Call2Go apresentaram média de 46,1 milhões de visualizações (mediana: 7,9M; DP: 98,8M) versus 19,1 milhões para vídeos sem Call2Go (mediana: 6,0M; DP: 44,0M). A escala logarítmica adotada no boxplot evidencia a distribuição altamente assimétrica à direita em ambos os grupos, com outliers expressivos correspondentes a vídeos virais de artistas como Marília Mendonça e Jorge \& Mateus. A diferença de médias (2,4$\times$) é relevante descritivamente, mas deve ser interpretada com cautela: artistas de maior popularidade tendem simultaneamente a ter mais views \textbf{e} a adotar práticas profissionalizadas de \textit{marketing} (incluindo Call2Go), o que configura potencial variável confundidora.

\subsection{Testes de Hipótese}

A Tabela~\ref{tab:estatisticos} sintetiza os resultados de todos os testes estatísticos conduzidos.

\begin{table}[H]
\centering
\caption{Resultados dos testes estatísticos para as hipóteses $H_2$, $H_3$ e análise bidirecional.}
\label{tab:estatisticos}
\renewcommand{\arraystretch}{1.3}
\begin{tabular}{@{}llccc@{}}
\toprule
\textbf{Hipótese / Análise} & \textbf{Teste} & \textbf{Estatística} & \textbf{\textit{p}-valor} & \textbf{Decisão} \\
\midrule
$H_2$: Call2Go $\to$ Views YT & Mann-Whitney U & $U = 118.004$ & $p = 0{,}360$ & Não rejeita $H_0$ \\
$H_3$: Call2Go $\to$ Pop. Spotify & Mann-Whitney U & $U = 169.415{,}5$ & $p = 0{,}998$ & Não rejeita $H_0$ \\
\midrule
Dir. A: Call2Go Rate $\leftrightarrow$ Pop. Spotify & Spearman ($\rho$) & $\rho = -0{,}086$ & $p = 0{,}556$ & Não significativo \\
Dir. A: Call2Go Rate $\leftrightarrow$ Seguidores & Spearman ($\rho$) & $\rho = -0{,}175$ & $p = 0{,}228$ & Não significativo \\
\midrule
Dir. B: Pop. Spotify $\leftrightarrow$ Avg Views YT & Spearman ($\rho$) & $\rho = 0{,}392$ & $p = 0{,}005$ & \textbf{Significativo ***} \\
Dir. B: Pop. Spotify $\leftrightarrow$ Avg Engagement & Spearman ($\rho$) & $\rho = 0{,}310$ & $p = 0{,}030$ & \textbf{Significativo **} \\
Dir. B: Seguidores $\leftrightarrow$ Avg Views YT & Spearman ($\rho$) & $\rho = 0{,}342$ & $p = 0{,}016$ & \textbf{Significativo **} \\
Dir. B: Seguidores $\leftrightarrow$ Avg Engagement & Spearman ($\rho$) & $\rho = 0{,}380$ & $p = 0{,}007$ & \textbf{Significativo ***} \\
Dir. B: Pop. $\leftrightarrow$ Views (por vídeo, $n=980$) & Spearman ($\rho$) & $\rho = 0{,}344$ & $p < 0{,}001$ & \textbf{Significativo ***} \\
\bottomrule
\multicolumn{5}{l}{\footnotesize *** $p < 0{,}01$; ** $p < 0{,}05$; $\alpha = 0{,}05$ para $H_2$ e $H_3$; $\alpha = 0{,}10$ para análise bidirecional (Spearman).}
\end{tabular}
\end{table}

Os resultados para $H_2$ ($p = 0{,}360$) e $H_3$ ($p = 0{,}998$) levam à \textbf{não-rejeição} da hipótese nula em ambos os casos. Para $H_2$, a interpretação é que a diferença de views observada descritivamente (46,1M vs.\ 19,1M) não é estatisticamente distinguível do esperado por acaso ao nível de 5\%, sugerindo que a variável confundidora (fama geral do artista) explicaria tanto a maior quantidade de views quanto a adoção de Call2Go, sem relação causal entre os dois. Para $H_3$, o quase-perfeito $p = 0{,}998$ indica que os grupos são praticamente indistinguíveis em termos de popularidade no Spotify (74,77 vs.\ 73,58 pontos em média), reforçando a ausência de evidência de impacto.

A análise bidirecional classifica a relação YouTube $\leftrightarrow$ Spotify como \textbf{UNIDIRECIONAL: Spotify $\to$ YouTube}. Todas as correlações da Direção B são estatisticamente significativas ($p \leq 0{,}030$), enquanto nenhuma das correlações da Direção A atinge significância ($p \geq 0{,}228$). Interpreta-se que artistas com maior presença consolidada no Spotify --- medida por \texttt{popularity} e \texttt{followers} --- tendem a acumular maior engajamento nos seus vídeos do YouTube, mas que a estratégia de Call2Go no YouTube não demonstra influência mensurável sobre as métricas do Spotify. É importante registrar que esta relação é correlacional, não causal: variáveis confundidoras não observadas (fama acumulada, investimento em marketing, qualidade do repertório) podem explicar parcialmente ambas as métricas simultaneamente.

\subsection{Validação Cruzada: Confiabilidade do Detector}

A Tabela~\ref{tab:crossval} apresenta as métricas de confiabilidade do detector em cada um dos três níveis de validação.

\begin{table}[H]
\centering
\caption{Métricas de confiabilidade do detector automatizado por nível de análise (amostra: 50 vídeos, \textit{ground truth} humano).}
\label{tab:crossval}
\renewcommand{\arraystretch}{1.3}
\begin{tabular}{@{}lcccccc@{}}
\toprule
\multirow{2}{*}{\textbf{Nível}} & \multirow{2}{*}{\textbf{Acurácia}} & \multicolumn{3}{c}{\textbf{\texttt{link\_direto}}} & \multicolumn{2}{c}{\textbf{\texttt{nenhum}}} \\
\cmidrule(lr){3-5} \cmidrule(lr){6-7}
 & & \textbf{P} & \textbf{R} & \textbf{F1} & \textbf{P} & \textbf{R} \\
\midrule
1 --- Só vídeo & 60,0\% & 100,0\% & 31,0\% & 47,4\% & 51,2\% & 100,0\% \\
\textbf{2 --- Combinado (vídeo + canal)} & \textbf{100,0\%} & \textbf{100,0\%} & \textbf{100,0\%} & \textbf{100,0\%} & \textbf{100,0\%} & \textbf{100,0\%} \\
3 --- Canal isolado & 100,0\% & 100,0\% & 100,0\% & 100,0\% & 100,0\% & 100,0\% \\
\bottomrule
\multicolumn{7}{l}{\footnotesize P = Precisão; R = \textit{Recall}; F1 = F1-Score; ($n_{\text{link\_direto}} = 29$; $n_{\text{nenhum}} = 21$).}
\end{tabular}
\end{table}

Os resultados da validação cruzada rejeitam $H_0$ com folga: o detector no Nível~2 atinge 100\% de acurácia, precisão e recall, superando amplamente o limiar de 85\% estabelecido em $H_1$. O padrão de erros no Nível~1 (Só vídeo) é revelador: a precisão de 100\% com recall de apenas 31\% evidencia que, quando o sistema detecta Call2Go no texto da descrição do vídeo, ele \textbf{nunca} erra (zero falsos positivos), mas deixa de detectar 69\% dos casos reais de Call2Go --- que estão presentes no perfil do canal, não na descrição individual do vídeo. Esse padrão confirma que a prática dominante dos artistas da amostra é inserir o link do Spotify no canal de forma global, e não em cada vídeo individualmente. A recuperação total ao Nível~3 (Canal isolado) indica que o \textit{scraping} da aba ``Sobre'' fornece evidências suficientes para a classificação correta de todos os 50 casos da amostra.

\subsection{Limitações e Direcionamentos Futuros}

Este trabalho apresenta limitações que devem ser explicitadas para adequada contextualização de seus resultados. Em primeiro lugar, a coleta de métricas do Spotify é \textit{snapshot} transversal (30/03/2026), sem seriação temporal, o que impede a análise de causalidade entre mudanças nas estratégias de Call2Go e variações posteriores na popularidade. A construção de uma série temporal exigiria coletas mensais repetidas, o que excede o escopo deste TCC mas constitui direcionamento natural para pesquisa futura. Em segundo lugar, a análise bidirecional da Direção B é inteiramente correlacional e não utiliza dados de links reais do Spotify apontando para o YouTube (que a Spotify Web API não expõe), de forma que a relação identificada pode decorrer de variável confundidora não observada representativa de fama ou notoriedade geral do artista. Em terceiro lugar, o corpus é limitado ao segmento de artistas brasileiros de maior notoriedade, o que restringe a generalização dos achados quantitativos para artistas independentes ou de nicho, embora a metodologia de detecção seja transferível a qualquer corpus.

% ============================================================
%  7. CONCLUSÃO
% ============================================================
\section{Conclusão}
\label{sec:conclusao}

Este trabalho apresentou o desenvolvimento e a validação de um sistema automatizado de detecção de estratégias Call2Go em vídeos do YouTube, aplicado ao contexto do mercado musical digital brasileiro. O principal resultado metodológico --- acurácia, precisão e recall de 100\% no nível de análise combinado (vídeo + canal) --- valida a abordagem baseada em expressões regulares multicamada como instrumento confiável para análise de conteúdo em larga escala, respondendo afirmativamente à hipótese $H_1$ e rejeita $H_0$. A decomposição da detecção em três níveis analíticos e a complementação da coleta via API com \textit{scraping} estruturado das informações de canal emergiram como decisões metodológicas determinantes para o alcance de acurácia plena.

No plano substantivo, os testes estatísticos não encontraram evidências significativas de que a adoção de Call2Go impacta o volume de visualizações no YouTube ($H_2$: $p = 0{,}360$) ou a popularidade no Spotify ($H_3$: $p = 0{,}998$), levando à não-rejeição da hipótese nula em ambos os casos. A análise bidirecional identificou uma relação unidirecional estatisticamente significativa no sentido Spotify $\to$ YouTube ($\rho = 0{,}392$; $p = 0{,}005$), sugerindo que artistas com maior presença no Spotify tendem a acumular mais engajamento no YouTube, sem que o inverso seja observado nos dados. Esses resultados não implicam que a estratégia Call2Go seja ineficaz, mas indicam que seus efeitos, se existentes, operam em escalas ou janelas temporais que a metodologia transversal deste estudo não consegue capturar.

As contribuições deste trabalho são de natureza primariamente metodológica: (i) a consolidação de um pipeline ETL reprodutível e documentado para coleta e análise de dados musicais cross-platform; (ii) o protocolo de validação cruzada em três níveis, aplicável a outros contextos de classificação automatizada de conteúdo textual; (iii) o detector Call2Go com documentação completa de seus padrões e filtros, pronto para reutilização em outros corpora; e (iv) a demonstração de que o \textit{scraping} estruturado das informações de canal é metodologicamente imprescindível para análise correta do fenômeno, sendo a omissão dessa fonte suficiente para reduzir a acurácia de 100\% para 60\%.

% ============================================================
%  REFERÊNCIAS
% ============================================================
\begin{thebibliography}{99}

\bibitem{ifpi2023}
IFPI. \textit{Global Music Report 2023}. London: International Federation of the Phonographic Industry, 2023. Disponível em: \url{https://www.ifpi.org/global-music-report/}. Acesso em: 28 mar. 2026.

\bibitem{mann1947}
MANN, H. B.; WHITNEY, D. R. On a test of whether one of two random variables is stochastically larger than the other. \textit{The Annals of Mathematical Statistics}, v.~18, n.~1, p.~50--60, 1947.

\bibitem{cohen1960}
COHEN, J. A coefficient of agreement for nominal scales. \textit{Educational and Psychological Measurement}, v.~20, n.~1, p.~37--46, 1960.

\bibitem{fleiss1971}
FLEISS, J. L. Measuring nominal scale agreement among many raters. \textit{Psychological Bulletin}, v.~76, n.~5, p.~378--382, 1971.

\bibitem{krippendorff2004}
KRIPPENDORFF, K. \textit{Content Analysis: An Introduction to Its Methodology}. 2.~ed. Thousand Oaks: Sage Publications, 2004.

\bibitem{schedl2014}
SCHEDL, M.; GÓMEZ, E.; URBANO, J. Music information retrieval: Recent developments and applications. \textit{Foundations and Trends in Information Retrieval}, v.~8, n.~2--3, p.~127--261, 2014.

\bibitem{celma2010}
CELMA, O. \textit{Music Recommendation and Discovery: The Long Tail, Long Fail, and Long Play in the Digital Music Space}. Berlin: Springer, 2010.

\bibitem{wlomert2016}
WLÖMERT, N.; PAPIES, D. On-demand streaming services and music industry revenues --- insights from Spotify's market entry. \textit{International Journal of Research in Marketing}, v.~33, n.~2, p.~314--327, 2016.

\bibitem{salganik2019}
SALGANIK, M. J. \textit{Bit by Bit: Social Research in the Digital Age}. Princeton: Princeton University Press, 2019. Disponível em: \url{https://www.bitbybitbook.com}. Acesso em: 15 mar. 2026.

\bibitem{aguiar2018}
AGUIAR, L.; WALDFOGEL, J. Platforms, power and promotion: Distribution dynamics and digital music. \textit{The Journal of Industrial Economics}, v.~66, n.~3, p.~615--653, 2018.

\bibitem{napoli2012}
NAPOLI, P. M. \textit{Audience Evolution: New Technologies and the Transformation of Media Audiences}. New York: Columbia University Press, 2012.

\end{thebibliography}

\end{document}
